\begin{proposition}[Refuted: even normal local domains of characteristic zero]
\label{C:normal_domain_counterexample}
Let
\[
 R=\left(\mathbb Z_{(2)}[a,b,c]/(bc-a(2-a))\right)_{(2,a,b,c)}.
\]
Then $R$ is a normal noetherian local domain of characteristic zero,
$K_0(R)=\mathbb Z$, and $R\otimes\mathbb Q=R[1/2]$ is a connected regular
ring. Nevertheless $\vartheta_R$ is not surjective in degree two.
\end{proposition}
\begin{proof}
The polynomial $bc-a(2-a)$ is primitive and irreducible as a degree-one
polynomial in $c$ over the UFD $\mathbb Z_{(2)}[a,b]$: its two coefficients
$b$ and $-a(2-a)$ are relatively prime. Thus $R$ is a domain of
characteristic zero. It is a hypersurface in a regular local ring, hence
Cohen--Macaulay and satisfies $(S_2)$. After inverting two its defining
hypersurface is smooth: the three partial derivatives are $2a-2,c,b$,
and their simultaneous vanishing would give $a=1,b=c=0$, inconsistent
with the equation. The ideal $(2)$ is prime, since its quotient is the
local domain $\mathbb F_2[a,b,c]/(a^2+bc)$. The quadratic polynomial has
no linear factor; this is also checked by the attached exact program.
Consequently the only height-one prime containing two is $(2)$ itself,
and $R_{(2)}$ is regular of dimension one with parameter two. Thus $R$
satisfies $(R_1)$ and is normal by Serre's criterion
\cite[Tag 031S]{C-StacksNormal}. The stated connectedness and regularity of
$B=R[1/2]$ follow at once, and locality gives $K_0(R)=\mathbb Z$.

Put $I=(a,b)$. This is a height-one prime: the quotient is
$\mathbb Z_{(2)}[c]_{(2,c)}$, of dimension two, while $\dim R=3$.
It is not principal. Indeed the hypersurface equation lies in the square
of the ambient maximal ideal, so the images of $2,a,b,c$ are linearly
independent in $\mathfrak m/\mathfrak m^2$. In particular $a,b$ remain
independent in $I/\mathfrak m I$, and Nakayama's lemma says that $I$ needs
two generators. In a normal domain a height-one prime is the intersection
of its localizations at height-one primes, inside the fraction field:
at $I$ this follows by contraction of the prime, and at every other
height-one prime its localization is the whole ring. The intersection
statement for $R$ is \cite[Tag 031T]{C-StacksNormal}.

We claim that $IB$ is still not principal. If $IB=xB$, let $v_{(2)}(x)=m$
and replace $x$ by $y=2^{-m}x$. For every height-one prime other than $(2)$,
localization has already inverted two and gives $I R_{\mathfrak p}=y
R_{\mathfrak p}$. At $(2)$, $I R_{(2)}=R_{(2)}$ because $a\notin(2)$,
and the choice of $y$ gives the same equality. Intersecting these equalities
at all height-one primes forces $I=yR$, a contradiction. This argument
requires no computation of the entire divisor class group.

Now let the transposition in $\mathfrak S_2$ act on $R^2$ by
\[
 T=\begin{pmatrix}a-1&b\\c&1-a\end{pmatrix}.
\]
Its square is the identity, since $T^2-1=(bc-a(2-a))1_2$.
Over $B$ the idempotent $e=(T+1)/2$ has rank one. Projection to the first
coordinate identifies $\operatorname{im}(e)$ with $IB$: its image is
$(a,b)B$, and its kernel is zero because it is a map from a rank-one
torsion-free module that is nonzero generically. Thus $IB$ is invertible
but nontrivial in $\operatorname{Pic}(B)$.

The exact functor taking the positive eigenspace after extension to $B$
induces $R_R(\mathfrak S_2)\to K_0(B)$; composing with determinant gives
an additive map to $\operatorname{Pic}(B)$. Both $\rho_2$ and $\rho_1^2$
have a free positive eigenspace, so their images are zero. The representation
$T$ has image $[IB]\ne0$. Hence its class is outside the degree-two image
of $\vartheta_R$, as claimed. The symbolic matrix identity and the
mod-two irreducibility check appear in
\texttt{code/C\_involution\_certificate.py} and its JSON output.
\end{proof}
